\section{Limitations and Future Work}
\label{sec:limitations}

A clear statement of limitations is part of any honest analysis.
This section lists the seven most important constraints and, where
relevant, the natural next step that would relax each one.

\subsection{Target substitution from T-Bill yield to WACMR}

The original project proposal targeted the 364-day Treasury Bill
cut-off yield. As documented in Section~\ref{sec:database}, the NDAP
\texttt{Treasury\_Bills\_Details.csv} table publishes outstanding
amounts rather than yields, despite the schema label suggesting
otherwise. WACMR was substituted because it is available, complete,
and arguably a superior target for studying short-rate transmission.
This is a substantively different target from the original proposal
and the report's claims should be read with that in mind.

\textbf{Future work:} If RBI adds T-Bill cut-off yields as a weekly
machine-readable series (Recommendation 4 in Section
\ref{sec:recommendations}), the same pipeline can be retrained with
the original target and the results compared.

\subsection{Two-regime simplicity}

The K-Means clustering at $K=2$ is statistically optimal by silhouette
score, but the silhouette difference between $K=2$ (0.4238) and
$K=3$ (0.4048) is only 0.0190, which is within noise. A three-regime
fit would split Regime 1 into a COVID-emergency sub-period and a
post-2022 normalisation sub-period, both of which are macroeconomically
meaningful. The two-regime choice was made on interpretability grounds:
the $K=3$ split does not align with any single observable RBI policy
boundary.

\textbf{Future work:} A hierarchical clustering (agglomerative or HDBSCAN)
might produce a more nuanced regime structure that does not require an
arbitrary $K$ choice.

\subsection{Two columns retained with substantial imputation}

\texttt{rates\_I7496\_28} (Market Repo Rate) and
\texttt{rates\_I7496\_29} (Certificate of Deposit Rate) had 17.4\%
missing values, marginally below the 75\% density threshold. They
were retained with median imputation because of their economic
relevance. A sensitivity analysis re-running the K-Means and the
walk-forward XGBoost without these two columns is worth doing; the
expected delta is small but not zero.

\subsection{Single-horizon forecasting only}

The walk-forward analysis predicts WACMR one week ahead. Practical
treasury planning often needs two-week and four-week horizons, and
operational liquidity decisions can need monthly or quarterly
forecasts. A multi-horizon forecasting setup, either as separate
models per horizon or as a recursive one-step-ahead generator, is the
natural extension.

\textbf{Future work:} Conformal prediction or quantile XGBoost would
also add proper uncertainty quantification rather than the current
implicit assumption of homoskedastic Gaussian residuals.

\subsection{Single-target regression rather than vector autoregression}

The forecasting setup treats WACMR as a single scalar target and
ignores the joint dynamics with related rates (91-day T-Bill yield,
10-year G-Sec yield, USD/INR). A vector autoregression or a
gradient-boosted multi-target regressor would capture the cross-rate
dependencies that drive the term structure of short rates.

\subsection{No causal identification claim}

The counterfactual simulator (Section~\ref{sec:system}) is a
model-based prediction, not a causal experiment. When the user moves
the Repo Rate slider, the simulator returns the prediction the model
would make if the Repo Rate were at that level, holding all other
features fixed at their current observed values. This is not the same
as the causal effect of an actual policy intervention, because real
policy changes shift many features simultaneously.

\textbf{Future work:} Combining the model with a structural VAR or
with synthetic-control methods could identify causal effects of
specific historical policy actions. The 2022 inter-meeting hike is a
natural candidate.

\subsection{Manual NLP corpus does not scale}

The seventy-five-event corpus in Section~\ref{sec:methods-nlp} is
expert-curated and produces clean labels, but it is impossible to
scale manually beyond a few hundred events. The evidence in
Section~\ref{sec:results} shows that the existing corpus correlation
with WACMR is real but weak ($r=-0.2081$); a denser corpus might or
might not surface a stronger relationship.

\textbf{Future work:} An automated FinBERT or LLM-based sentiment
pipeline over the full RBI press-release archive (about 5,000
documents over the 2014--2024 window) is the obvious extension.
Domain fine-tuning on a small labelled set of Indian-monetary-policy
language would address the off-the-shelf model limitations identified
in Section~\ref{sec:methods-nlp}.

\subsection{No event dummies for extraordinary policy actions}

The residual analysis (Section~\ref{sec:methods-supervised}) shows
that the model's tail errors concentrate on three or four
unsignalled inter-meeting actions. Including binary event dummies
for these actions, derived from the published MPC calendar and from
out-of-cycle press releases, is the obvious remediation. It was not
implemented in this project because it requires a separate data
source and the marginal cost was higher than the marginal benefit
within the project timeline.
